%%%%%%%%%%%%%%%%%%%%%%%%%%%%%%%%%%%%%%%%%%%%%%%%%%%%%%%%%%%%%%%%%%%%%%%%%%%%%%%
%%%%%%%%%%%%%%%%%%%%%%%%%%%%%%%%%%%%%%%%%%%%%%%%%%%%%%%%%%%%%%%%%%%%%%%%%%%%%%%
%%%%%%%%%%%%%%%%%%%%%%%%%%%%%%%%%%%%%%%%%%%%%%%%%%%%%%%%%%%%%%%%%%%%%%%%%%%%%%%

% Typeset by LaTeX by Lubomir Stepanek, September 2017.


%%%%%%%%%%%%%%%%%%%%%%%%%%%%%%%%%%%%%%%%%%%%%%%%%%%%%%%%%%%%%%%%%%%%%%%%%%%%%%%

%% I am defining class and properties of the document -------------------------
\documentclass[t, compress]{beamer}


%% ----------------------------------------------------------------------------
%% I am setting title page and global parameters ------------------------------
%% ----------------------------------------------------------------------------

\mode<presentation>{

%% I am including my own styles saved in "my_styles.text" file ----------------
%%%%%%%%%%%%%%%%%%%%%%%%%%%%%%%%%%%%%%%%%%%%%%%%%%%%%%%%%%%%%%%%%%%%%%%%%%%%%%%
%%%%%%%%%%%%%%%%%%%%%%%%%%%%%%%%%%%%%%%%%%%%%%%%%%%%%%%%%%%%%%%%%%%%%%%%%%%%%%%
%%%%%%%%%%%%%%%%%%%%%%%%%%%%%%%%%%%%%%%%%%%%%%%%%%%%%%%%%%%%%%%%%%%%%%%%%%%%%%%

% presentation theme ----------------------------------------------------------
\usetheme{Frankfurt}

% mathematical notation of article kind ---------------------------------------
\usefonttheme[onlymath]{serif}

% I am loading some packages --------------------------------------------------
\usepackage[T1]{fontenc}
\usepackage[utf8]{inputenc}
\usepackage[german]{babel}
\usepackage{csquotes}
\usepackage{amsmath}
\usepackage{physics}
\usepackage{amsfonts}
\usepackage{array}
\usepackage{multirow}
\usepackage{graphicx}
\usepackage{epstopdf}   % it guarantees .eps figures are displayed
                        % properly on Windows
\usepackage{blindtext}
\usepackage{scrextend}
\usepackage{remreset}
\usepackage{hyperref}
\usepackage{listings}
\usepackage{xcolor}
\usepackage{subfig}
\usepackage{tikz}
\usetikzlibrary{arrows,positioning} 
\tikzset{
    %Define standard arrow tip
    >=stealth',
    %Define style for boxes
    punkt/.style={
           rectangle,
           rounded corners,
           draw=black, very thick,
           text width=6.5em,
           minimum height=2em,
           text centered},
    % Define arrow style
    pil/.style={
           ->,
           thick,
           shorten <=2pt,
           shorten >=2pt,}
}

% I am setting a frame size and margins ---------------------------------------
\setbeamersize{%
  text margin left = 0.5cm,
  text margin right = 0.5cm
}

% new command for reduced font size -------------------------------------------
\newcommand\reduceFontSize{\fontsize{8}{10}\selectfont}

% I am setting a graphic path -------------------------------------------------
\graphicspath{{./figures/}}

% I am defining my own footer that is unique for title page -------------------
\defbeamertemplate*{footline}{my footline}
{
    \ifnum \insertpagenumber = 1
        \leavevmode%
        \hbox{%
            \pgfsetfillopacity{1}\begin{beamercolorbox}[wd=.20\paperwidth,
              ht=2.25ex,dp=1ex,center]{author in head/foot}%
              \usebeamerfont{author in head/foot}\pgfsetfillopacity{1}
            \end{beamercolorbox}%
            \pgfsetfillopacity{1}\begin{beamercolorbox}[wd=.60\paperwidth,
              ht=2.25ex,dp=1ex,center]{title in head/foot}%
              \usebeamerfont{title in head/foot}\pgfsetfillopacity{1}
            \end{beamercolorbox}%
            \pgfsetfillopacity{.8}\begin{beamercolorbox}[wd=.14\paperwidth,
              ht=2.25ex,dp=1ex,center]{date in head/foot}%
              \usebeamerfont{date in head/foot}\pgfsetfillopacity{1}%
              \insertshortdate{}
            \end{beamercolorbox}%
            \pgfsetfillopacity{1}\begin{beamercolorbox}[wd=.06\paperwidth,
              ht=2.25ex,dp=1ex,right]{date in head/foot}%
              \insertframenumber{}/\inserttotalframenumber\hspace*{1ex}
            \end{beamercolorbox}
          }%
        \vskip0pt%
    \else
        \leavevmode%
        \hbox{%
            \pgfsetfillopacity{1}\begin{beamercolorbox}[wd=.20\paperwidth,
              ht=2.25ex,dp=1ex,center]{author in head/foot}%
              \usebeamerfont{author in head/foot}\pgfsetfillopacity{1}%
                \hspace*{0.5ex}\insertshortauthor
            \end{beamercolorbox}%
            \pgfsetfillopacity{1}\begin{beamercolorbox}[wd=.60\paperwidth,
              ht=2.25ex,dp=1ex,center]{title in head/foot}%
              \usebeamerfont{title in head/foot}\pgfsetfillopacity{1}%
                \insertshorttitle
            \end{beamercolorbox}%
            \pgfsetfillopacity{.8}\begin{beamercolorbox}[wd=.14\paperwidth,
              ht=2.25ex,dp=1ex,center]{date in head/foot}%
              \usebeamerfont{date in head/foot}\pgfsetfillopacity{1}%
                \insertshortdate{}
            \end{beamercolorbox}%
            \pgfsetfillopacity{1}\begin{beamercolorbox}[wd=.06\paperwidth,
              ht=2.25ex,dp=1ex,right]{date in head/foot}%  
              \insertframenumber{}/\inserttotalframenumber\hspace*{1ex}
            \end{beamercolorbox}
          }%
        \vskip0pt%
    \fi
}

\setbeamertemplate{footline}[my footline]

% navigation circles in header ------------------------------------------------
\makeatletter
    \@removefromreset{subsection}{section}
\makeatother
\setcounter{subsection}{1}

% I am defining my own commands to switch-off including of some frames
% into navigational circles ---------------------------------------------------
\makeatletter
\let\beamer@writeslidentry@miniframeson = \beamer@writeslidentry
\def\beamer@writeslidentry@miniframesoff{%
  \expandafter\beamer@ifempty\expandafter{\beamer@framestartpage}{}
  % does not happen normally
  {%else
    % removed \addtocontents commands
    \clearpage\beamer@notesactions%
  }
}
\newcommand*{\miniframeson}{%
  \let\beamer@writeslidentry = \beamer@writeslidentry@miniframeson
}
\newcommand*{\miniframesoff}{%
  \let\beamer@writeslidentry = \beamer@writeslidentry@miniframesoff
}

% removing a section from header and outline via "\section*" command ----------
\long\def\beamer@@ssection*#1{\beamer@section[]{}}

\makeatother

% I am defining highlighting of R language code chunks ------------------------
%\captionsetup[lstlisting]{position = bottom}
\definecolor{my_gray}{rgb}{0.4,0.4,0.4}
\renewcommand{\lstlistingname}{Code}  %% code label redefining
\lstdefinestyle{custom_R}{
  basicstyle = \ttfamily\scriptsize,
  belowcaptionskip = 1\baselineskip,
  breaklines = true,
  frame = L,
  xleftmargin = \parindent,
  language = R,
  showstringspaces = false,
  keywordstyle = \bfseries\color{green!40!black},
  commentstyle = \itshape\color{purple!40!black},
  identifierstyle = \color{blue},
  stringstyle = \color{orange},
  numbers = left,
  numbersep = 12pt,
  numberstyle = \small\color{my_gray},
  texcl = true,
  alsoother = {\#},
  inputencoding = utf8,
  extendedchars = true,
  literate = %
  {á}{{\'a}}1
  {č}{{\v{c}}}1
  {ď}{{\v{d}}}1
  {é}{{\'e}}1
  {ě}{{\v{e}}}1
  {í}{{\'i}}1
  {ň}{{\v{n}}}1
  {ó}{{\'o}}1
  {ř}{{\v{r}}}1
  {š}{{\v{s}}}1
  {ť}{{\v{t}}}1
  {ú}{{\'u}}1
  {ů}{{\r{u}}}1
  {ý}{{\'y}}1
  {ž}{{\v{z}}}1
  {Á}{{\'A}}1
  {Č}{{\v{C}}}1
  {Ď}{{\v{D}}}1
  {É}{{\'E}}1
  {Ě}{{\v{E}}}1
  {Í}{{\'I}}1
  {Ň}{{\v{N}}}1
  {Ó}{{\'O}}1
  {Ř}{{\v{R}}}1
  {Š}{{\v{S}}}1
  {Ť}{{\v{T}}}1
  {Ú}{{\'U}}1
  {Ů}{{\r{U}}}1
  {Ý}{{\'Y}}1
  {Ž}{{\v{Z}}}1
}
\lstset{escapechar = @, style = custom_R}

% I am defining a font of labels in labeling environtment ---------------------
\addtokomafont{labelinglabel}{\ttfamily}

% I am defining czech type of double quotes -----------------------------------
\DeclareQuoteAlias{german}{czech}
\MakeOuterQuote{"}


%%%%%%%%%%%%%%%%%%%%%%%%%%%%%%%%%%%%%%%%%%%%%%%%%%%%%%%%%%%%%%%%%%%%%%%%%%%%%%%
%%%%%%%%%%%%%%%%%%%%%%%%%%%%%%%%%%%%%%%%%%%%%%%%%%%%%%%%%%%%%%%%%%%%%%%%%%%%%%%
%%%%%%%%%%%%%%%%%%%%%%%%%%%%%%%%%%%%%%%%%%%%%%%%%%%%%%%%%%%%%%%%%%%%%%%%%%%%%%%







}


%% I am setting title page ----------------------------------------------------

\title[%
  Aplikace pro podporu správného výběru
  a realizace statistické metody%
]{%
  Aplikace \texttt{statisticke\_nastroje} \\
  pro podporu \\ správného výběru
  a realizace statistické metody%
}

\subtitle{%
  \vspace{-0.15cm}
  --- \\
  \vspace{-0.05cm}
  Letní škola matematické biologie 2017%
}

\author[%
  Lubomír Štěpánek%
]{%
  \vspace{-0.15cm}Lubomír Štěpánek\textsuperscript{1, 2}\vspace{-0.15cm}
}

\institute{%

  \begin{minipage}{.30\textwidth}
    \vspace{+0.15cm}
    \centering
    {%
      \hspace{+1.5cm}
      \includegraphics[width = 1.2cm]{lf1_logo.eps} \\
      \vspace{0.35cm}
      \hspace{+1.5cm}
      \includegraphics[width = 1.1cm]{fbmi_logo.eps}%
    }
  \end{minipage}%
  %
  \begin{minipage}{.70\textwidth}
    \textsuperscript{1}BIOSTAT \& Oddělení výpočetní techniky \\
    \phantom{\textsuperscript{1}}Ústav biofyziky a informatiky \\
    \phantom{\textsuperscript{1}}1. lékařská fakulta \\
    \phantom{\textsuperscript{1}}Univerzita Karlova v Praze \\
    \\
    \textsuperscript{2}Katedra biomedicínské informatiky \\
    \phantom{\textsuperscript{2}}Fakulta biomedicínského inženýrství \\
    \phantom{\textsuperscript{2}}České vysoké učení technické v Praze
  \end{minipage}
  
}

\date[14. září 2017]{\phantom{14. září 2017}}


%% ----------------------------------------------------------------------------
%% I am beginning the document ------------------------------------------------
%% ----------------------------------------------------------------------------

\begin{document}


%% title page -----------------------------------------------------------------

{
\setbeamertemplate{headline}{}
\miniframesoff % it is switching off including of the frame
               % in header navigational circles
\begin{frame}
  \titlepage
\end{frame}

\miniframeson  % it is switching on the including of the frame
               % in the header navigational circles again
}


%% talk content ---------------------------------------------------------------

\miniframesoff
\begin{frame}
  \frametitle{%
    Obsah ústního sdělení%
  }
  \tableofcontents
\end{frame}
\miniframeson


%% ----------------------------------------------------------------------------
%% sekce "úvod" ---------------------------------------------------------------
%% ----------------------------------------------------------------------------

\section{Úvod}

%% rychlý úvod ----------------------------------------------------------------

\begin{frame}
  \frametitle{%
    Rychlý úvod%
  }
  \begin{itemize}
    \item statistika ve výzkumu hraje vždy důležitou roli
    \item různé statistické nástroje mají různé nevýhody
    \item webová customizovaná aplikace se snaží pokrýt výhody
    \begin{itemize}
      \item komerčních nástrojů
      \item volně dostupných nástrojů
    \end{itemize}
  \end{itemize}
\end{frame}


%% odkaz na aplikaci ----------------------------------------------------------

\begin{frame}
  \frametitle{%
    Odkaz na aplikaci%
  }
  \vspace{+2.0cm}
  \begin{block}{\centering Odkaz}
    \center
    \href{%
      http://shiny.statest.cz:3838/statisticke\_nastroje/%
    }{%
      http://shiny.statest.cz:3838/statisticke\_nastroje/%
    }
  \end{block}
\end{frame}


%% ----------------------------------------------------------------------------
%% sekce "o aplikaci" ---------------------------------------------------------
%% ----------------------------------------------------------------------------

\section{O aplikaci}

%% prakticky o aplikaci -------------------------------------------------------

\begin{frame}
  \frametitle{%
    Prakticky o aplikaci%
  }
  \begin{itemize}
    \item cílí hlavně na česky mluvící pregraduální studenty
    \item podporuje fázi výběru a realizace pomocí vhodné statistické
    metody
    \item online, nevyžaduje instalaci či licenci    
  \end{itemize}
\end{frame}


%% technicky o aplikaci -------------------------------------------------------

\begin{frame}
  \frametitle{%
    Technicky o aplikaci%
  }
  \begin{itemize}
    \item napsána hlavně pomocí jazyka R a jeho knihovny \texttt{shiny}
    \item umístěna na R-kovém serveru 1. LF UK
    \item obsahuje cca 1400 kódujících řádků v jazyce R
    \item HTML, CSS a \texttt{javascript} zajišťují uživatelský komfort
  \end{itemize}
\end{frame}


%% ----------------------------------------------------------------------------
%% sekce "vývojové diagramy pro podporu volby vhodné metody" ------------------
%% ----------------------------------------------------------------------------

\section{Diagramy pro podporu volby vhodné metody}

%% vývojové diagramy pro pro podporu volby vhodné metody ----------------------

\begin{frame}
  \frametitle{}
  \begin{figure}
    \centering
    \vspace{-0.35cm}
    \includegraphics[width = 11.1cm]{%
      flow_chart.pdf%
    }
  \end{figure}
\end{frame}


%% vývojové diagramy pro pro podporu volby vhodné metody ----------------------

\begin{frame}
  \frametitle{}
  \begin{figure}
    \centering
    \vspace{-0.2cm}
    \includegraphics[width = 11.5cm]{%
      nice_flow_chart.pdf%
    }
  \end{figure}
\end{frame}


%% ----------------------------------------------------------------------------
%% sekce "aktuální stav" ------------------------------------------------------
%% ----------------------------------------------------------------------------

\section{Aktuální stav}

%% dosažené výsledky ----------------------------------------------------------

\begin{frame}
  \frametitle{%
    Dosažené výsledky%
  }
  \begin{block}{\centering Odkaz}
    \center
    \href{%
      http://shiny.statest.cz:3838/statisticke\_nastroje/%
    }{%
      http://shiny.statest.cz:3838/statisticke\_nastroje/%
    }
  \end{block}
  \begin{itemize}
    \item aplikace pokrývá běžné statistické metody
    \begin{itemize}
      \item parametrické metody
      \item neparametrické metody
    \end{itemize}
    \item implementován responzivní design
    \item v plánu zapojení aplikace pro účely analýz do závěrečných
    prací na \mbox{1. LF UK} v akademickém roce 2017/2018
  \end{itemize}
\end{frame}


%% uživatelské rozhraní -------------------------------------------------------

\begin{frame}
  \frametitle{%
    Uživatelské rozhraní%
  }  
  \vspace{-0.8cm}  
  \begin{figure}
    %\hspace{-1.0cm}
    \includegraphics[width = 12.0cm]{%
      user_interface.jpg%
    }
  \end{figure}
\end{frame}


%% ----------------------------------------------------------------------------
%% sekce "závěr" --------------------------------------------------------------
%% ----------------------------------------------------------------------------

\section{Závěrem}

%% conclusion -----------------------------------------------------------------

\begin{frame}
  \frametitle{%
    Shrnutí%
  }
  \begin{itemize}
    \item relativně krátkým kódem psaným v R získány pokročilé
    výpočetní metody volně dostupné online
    \item možnost "doimplementování" nových metod
    \item pomocí R je možné vytvářet i mnohozáložkové aplikace
    \mbox{s responzivním} designem
    \item serverové open-source R řešení plně podporuje češtinu
  \end{itemize}
\end{frame}


%% ----------------------------------------------------------------------------
%% section "díky za pozornost" ------------------------------------------------
%% ----------------------------------------------------------------------------

\section*{End}

%% děkuji za pozornost --------------------------------------------------------

{
\setbeamertemplate{headline}{}
\miniframesoff
\begin{frame}
\frametitle{}
  \vspace{+1.6cm}
  \begin{block}{\centering Děkuji za pozornost!}
    \center
    \begin{tabular}{l@{}l@{}l}
      &\href{mailto:lubomir.stepanek@lf1.cuni.cz}{%
        lubomir.stepanek@lf1.cuni.cz%
      } \\
      &\href{mailto:lubomir.stepanek@fbmi.cvut.cz}{%
        lubomir.stepanek@fbmi.cvut.cz%
      }
    \end{tabular}
  \end{block}
  \vspace{+0.8cm}
  \begin{center}
    
    \href{http://shiny.statest.cz:3838/statisticke\_nastroje/}{%
      \framebox{http://shiny.statest.cz:3838/statisticke\_nastroje/}%
    }
    
  \end{center}
\end{frame}
\miniframeson
}


%% ----------------------------------------------------------------------------
%% I am ending the document ---------------------------------------------------
%% ----------------------------------------------------------------------------

\end{document}


%%%%%%%%%%%%%%%%%%%%%%%%%%%%%%%%%%%%%%%%%%%%%%%%%%%%%%%%%%%%%%%%%%%%%%%%%%%%%%%
%%%%%%%%%%%%%%%%%%%%%%%%%%%%%%%%%%%%%%%%%%%%%%%%%%%%%%%%%%%%%%%%%%%%%%%%%%%%%%%
%%%%%%%%%%%%%%%%%%%%%%%%%%%%%%%%%%%%%%%%%%%%%%%%%%%%%%%%%%%%%%%%%%%%%%%%%%%%%%%





